% META: {"id": "1SPE-PROBA-COND-QCM", "chapitre": "1SPE-PROBA-COND", "type_objet": "qcm", "genere_depuis": "chapitres/1SPE-PROBA-COND/qcm/1SPE-PROBCOND-QCM.json", "status": "generated"}
% Fichier genere par scripts/build_qcm_tex.py — ne pas editer a la main.

\section*{\textcolor{chapcolor}{\MakeUppercase{Probabilites conditionnelles et independance --- QCM d'auto-évaluation}}}

\begin{center}
\textit{Pour chaque question, une seule réponse est exacte.}
\end{center}

\begin{enumerate}

\item[] \textbf{Capacité C1}

\item \textbf{[Q1]} On donne P(A) = 1/2 et P(A inter B) = 1/6. Que vaut P\_A(B) ?
  \begin{enumerate}[label=\Alph*.]
    \item 1/3
    \item 1/2
    \item 1/6
    \item 2/3
  \end{enumerate}

\item \textbf{[Q2]} P\_A(B) = 3/5 signifie :
  \begin{enumerate}[label=\Alph*.]
    \item P(A sachant B) = 3/5
    \item P(A inter B) = 3/5
    \item P(B sachant A) = 3/5
    \item P(A union B) = 3/5
  \end{enumerate}

\item \textbf{[Q3]} Si P(A) = 2/5 et P\_A(B) = 3/4, que vaut P(A inter B) ?
  \begin{enumerate}[label=\Alph*.]
    \item 3/20
    \item 3/10
    \item 23/20
    \item 6/5
  \end{enumerate}

\item \textbf{[Q4]} On donne P\_A(B) = 2/3. Que vaut P\_A(Bbar) ?
  \begin{enumerate}[label=\Alph*.]
    \item 2/3
    \item 1/3
    \item On ne peut pas savoir.
    \item 0
  \end{enumerate}

\item[] \textbf{Capacité C2}

\item \textbf{[Q5]} Dans un arbre pondere, la probabilite d'un chemin s'obtient en :
  \begin{enumerate}[label=\Alph*.]
    \item additionnant
    \item multipliant
    \item prenant le max
    \item divisant
  \end{enumerate}

\item \textbf{[Q6]} A chaque noeud d'un arbre pondere, la somme des probabilites des branches vaut :
  \begin{enumerate}[label=\Alph*.]
    \item 0
    \item la probabilite du noeud parent
    \item 1
    \item le nombre de branches
  \end{enumerate}

\item \textbf{[Q7]} Les probabilites inscrites sur les branches issues d'un noeud A sont :
  \begin{enumerate}[label=\Alph*.]
    \item P(B), P(Bbar)
    \item P\_A(B), P\_A(Bbar)
    \item P(A inter B), P(A inter Bbar)
    \item des frequences
  \end{enumerate}

\item[] \textbf{Capacité C3}

\item \textbf{[Q8]} Si P(B) = 1/3, P\_B(A) = 1/2 et P\_Bbar(A) = 1/4, alors P(A) vaut :
  \begin{enumerate}[label=\Alph*.]
    \item $1/4$
    \item $1/6$
    \item $5/12$
    \item $1/3$
  \end{enumerate}

\item \textbf{[Q9]} Pour appliquer la formule des probabilites totales avec {B1, B2, B3}, il faut que :
  \begin{enumerate}[label=\Alph*.]
    \item Les Bi soient independants
    \item Les Bi forment une partition
    \item Les Bi soient equiprobables
    \item P(Bi) > 1/2
  \end{enumerate}

\item \textbf{[Q10]} La formule des probabilites totales pour {B, Bbar} s'ecrit :
  \begin{enumerate}[label=\Alph*.]
    \item P(A) = P(B) + P(Bbar)
    \item P(A) = P(A inter B) x P(A inter Bbar)
    \item P(A) = P(B)*P\_B(A) + P(Bbar)*P\_Bbar(A)
    \item P(A) = P(B)*P(A) + P(Bbar)*P(A)
  \end{enumerate}

\item \textbf{[Q11]} Quelle etape fait partie de la demonstration de la formule des probabilites totales ?
  \begin{enumerate}[label=\Alph*.]
    \item A = A union B
    \item A = (A inter B) union (A inter Bbar) incompatibles
    \item P(A) = P(B) x P(A)
    \item A et B sont independants
  \end{enumerate}

\item[] \textbf{Capacité C4}

\item \textbf{[Q12]} A et B sont independants ssi :
  \begin{enumerate}[label=\Alph*.]
    \item P(A inter B) = 0
    \item P(A inter B) = P(A) + P(B)
    \item P(A inter B) = P(A) x P(B)
    \item P(A union B) = P(A) x P(B)
  \end{enumerate}

\item \textbf{[Q13]} Si A et B sont incompatibles avec P(A) > 0 et P(B) > 0, alors :
  \begin{enumerate}[label=\Alph*.]
    \item ils sont independants
    \item ils ne sont pas independants
    \item on ne peut pas conclure
    \item P\_A(B) = 1
  \end{enumerate}

\item \textbf{[Q14]} Si A et B sont independants, alors Abar et B sont :
  \begin{enumerate}[label=\Alph*.]
    \item incompatibles
    \item independants
    \item de meme probabilite
    \item sans lien
  \end{enumerate}

\item[] \textbf{Capacité C5}

\item \textbf{[Q15]} La sensibilite d'un test est :
  \begin{enumerate}[label=\Alph*.]
    \item P(M | T+)
    \item P(T+ | M)
    \item P(T- | Mbar)
    \item P(M)
  \end{enumerate}

\item \textbf{[Q16]} La valeur predictive positive (VPP) est :
  \begin{enumerate}[label=\Alph*.]
    \item P(T+ | M)
    \item P(M | T+)
    \item P(T+)
    \item P(Mbar | T-)
  \end{enumerate}

\item \textbf{[Q17]} Sensibilite 99\%, specificite 99\%, prevalence 0,1\%. La VPP vaut environ :
  \begin{enumerate}[label=\Alph*.]
    \item $99\%$
    \item $50\%$
    \item $9\%$
    \item $0{,}1\%$
  \end{enumerate}

\item \textbf{[Q18]} Un atelier connait les parts de production de trois machines et, pour chacune, la probabilite de produire une piece defectueuse. Pour determiner la machine d'origine sachant qu'une piece est defectueuse, on utilise :
  \begin{enumerate}[label=\Alph*.]
    \item l'independance
    \item formule de Bayes (prob totales + cond inversee)
    \item uniquement des prob simples
    \item theoreme des valeurs intermediaires
  \end{enumerate}

\end{enumerate}

% NEXUS-QCM-TEACHER-ONLY-BEGIN
\ifnxVersionProfesseur
\clearpage
\section*{Clé de correction — réservée au professeur}

\begin{center}
\begin{tabular}{lll}
\hline
Question & Capacité & Réponse exacte \\
\hline
Q1 & C1 & \textbf{A} \\
Q2 & C1 & \textbf{C} \\
Q3 & C1 & \textbf{B} \\
Q4 & C1 & \textbf{B} \\
Q5 & C2 & \textbf{B} \\
Q6 & C2 & \textbf{C} \\
Q7 & C2 & \textbf{B} \\
Q8 & C3 & \textbf{D} \\
Q9 & C3 & \textbf{B} \\
Q10 & C3 & \textbf{C} \\
Q11 & C3 & \textbf{B} \\
Q12 & C4 & \textbf{C} \\
Q13 & C4 & \textbf{B} \\
Q14 & C4 & \textbf{B} \\
Q15 & C5 & \textbf{B} \\
Q16 & C5 & \textbf{B} \\
Q17 & C5 & \textbf{C} \\
Q18 & C5 & \textbf{B} \\
\hline
\end{tabular}
\end{center}

\subsection*{Diagnostic des réponses erronées}

\begin{itemize}
  \item \textbf{Q1}
  \begin{itemize}
    \item \textbf{B} — L'eleve a confondu P\_A(B) avec P(A). P\_A(B) = P(A inter B)/P(A) = (1/6)/(1/2) = 1/3. \emph{Renvoi : C1, definition.}
    \item \textbf{C} — L'eleve a donne P(A inter B) au lieu de calculer le quotient. \emph{Renvoi : C1, formule.}
    \item \textbf{D} — L'eleve a calcule le complementaire de $P_A(B)=1/3$ et obtenu $1-1/3=2/3$ ; ce $2/3$ vaut $P_A(\overline B)$, pas $P_A(B)$. \emph{Renvoi : C1, formule.}
  \end{itemize}
  \item \textbf{Q2}
  \begin{itemize}
    \item \textbf{A} — L'eleve a inverse l'ordre de conditionnement : P\_A(B) est la probabilite de B sachant A, pas de A sachant B. \emph{Renvoi : C1, notation.}
    \item \textbf{B} — P\_A(B) n'est pas P(A inter B), c'est le quotient P(A inter B)/P(A). \emph{Renvoi : C1, definition.}
    \item \textbf{D} — Confusion entre probabilite conditionnelle et probabilite de l'union. \emph{Renvoi : C1, definition.}
  \end{itemize}
  \item \textbf{Q3}
  \begin{itemize}
    \item \textbf{A} — L'eleve a multiplie $1/5$ par $3/4$ en oubliant le facteur $2$ au numerateur de $P(A)=2/5$, ce qui donne exactement $3/20$ au lieu de $3/10$. \emph{Renvoi : C1, formule de multiplication.}
    \item \textbf{C} — L'eleve a additionne P(A) + P\_A(B) au lieu de les multiplier. \emph{Renvoi : C1, formule de multiplication.}
    \item \textbf{D} — L'eleve a multiplie les numerateurs $2 \times 3$ mais a oublie le denominateur $4$ de $P_A(B)=3/4$, obtenant $6/5$ au lieu de $(2 \times 3)/(5 \times 4)=3/10$. \emph{Renvoi : C1, formule de multiplication.}
  \end{itemize}
  \item \textbf{Q4}
  \begin{itemize}
    \item \textbf{A} — L'eleve a recopie P\_A(B) au lieu de calculer le complementaire. \emph{Renvoi : C1, complementaire conditionnel.}
    \item \textbf{C} — La relation P\_A(B) + P\_A(Bbar) = 1 est toujours vraie, donc on peut conclure. \emph{Renvoi : C1, propriete.}
    \item \textbf{D} — Confusion avec l'incompatibilite. \emph{Renvoi : C1, complementaire.}
  \end{itemize}
  \item \textbf{Q5}
  \begin{itemize}
    \item \textbf{A} — On additionne les probabilites des chemins pour obtenir la probabilite d'un evenement, mais le long d'un chemin on multiplie. \emph{Renvoi : C2, regle du produit.}
    \item \textbf{C} — Erreur de methode : on multiplie les probabilites le long d'un chemin. \emph{Renvoi : C2, regle du produit.}
    \item \textbf{D} — La division intervient dans le calcul de P\_A(B), pas dans le produit le long d'un chemin. \emph{Renvoi : C2, regle du produit.}
  \end{itemize}
  \item \textbf{Q6}
  \begin{itemize}
    \item \textbf{A} — Les probabilites sont positives, leur somme ne peut pas etre nulle. \emph{Renvoi : C2, regle de construction.}
    \item \textbf{B} — La somme des probabilites conditionnelles issues d'un noeud vaut 1, pas la probabilite du noeud. \emph{Renvoi : C2, regle de construction.}
    \item \textbf{D} — Confusion entre le nombre de branches et la somme des probabilites. \emph{Renvoi : C2, regle de construction.}
  \end{itemize}
  \item \textbf{Q7}
  \begin{itemize}
    \item \textbf{A} — Ce sont des probabilites conditionnelles, pas des probabilites marginales. \emph{Renvoi : C2, construction.}
    \item \textbf{C} — Les probabilites d'intersection sont celles des feuilles (chemins complets), pas des branches. \emph{Renvoi : C2, construction.}
    \item \textbf{D} — Un arbre pondere utilise des probabilites, pas des frequences observees. \emph{Renvoi : C2, definition.}
  \end{itemize}
  \item \textbf{Q8}
  \begin{itemize}
    \item \textbf{A} — L'eleve a recopie $P_{\overline B}(A)=1/4$ sans ponderer les deux branches de la partition. \emph{Renvoi : C3, formule.}
    \item \textbf{B} — L'eleve a calcule seulement $P(B)P_B(A)=1/6$ et oublie le second chemin $P(\overline B)P_{\overline B}(A)=1/6$. \emph{Renvoi : C3, formule.}
    \item \textbf{C} — L'eleve additionne $1/6$ et la probabilite conditionnelle $1/4$ sans la ponderer par $P(\overline B)=2/3$ ; le second terme vaut $1/6$. \emph{Renvoi : C3, formule.}
  \end{itemize}
  \item \textbf{Q9}
  \begin{itemize}
    \item \textbf{A} — L'independance n'est pas requise ; c'est la structure de partition qui est necessaire. \emph{Renvoi : C3, hypothese.}
    \item \textbf{C} — L'equiprobabilite n'est pas necessaire. \emph{Renvoi : C3, hypothese.}
    \item \textbf{D} — Si les trois evenements verifiaient chacun $P(B_i)>1/2$, la somme de leurs probabilites serait strictement superieure a 1 : ils ne pourraient donc pas former une partition. \emph{Renvoi : C3, hypothese.}
  \end{itemize}
  \item \textbf{Q10}
  \begin{itemize}
    \item \textbf{A} — P(B) + P(Bbar) = 1, pas P(A). Il manque les probabilites conditionnelles. \emph{Renvoi : C3, formule.}
    \item \textbf{B} — On additionne les termes, on ne les multiplie pas. \emph{Renvoi : C3, formule.}
    \item \textbf{D} — La formule D se simplifie en P(A) = P(A), ce qui est trivial et inutile. \emph{Renvoi : C3, formule.}
  \end{itemize}
  \item \textbf{Q11}
  \begin{itemize}
    \item \textbf{A} — A = A union B est faux en general (sauf si A inclus dans B). \emph{Renvoi : C3, demonstration.}
    \item \textbf{C} — Cela donnerait P(B) = 1, ce qui est faux en general. \emph{Renvoi : C3, demonstration.}
    \item \textbf{D} — L'independance n'est pas une hypothese de la demonstration. \emph{Renvoi : C3, demonstration.}
  \end{itemize}
  \item \textbf{Q12}
  \begin{itemize}
    \item \textbf{A} — P(A inter B) = 0 signifie incompatibilite, pas independance. \emph{Renvoi : C4, definition.}
    \item \textbf{B} — L'eleve a additionne les probabilites au lieu de les multiplier. Cette egalite viole aussi la borne $P(A\cap B)\leqslant \min(P(A),P(B))$, sauf dans le cas trivial ou les deux probabilites sont nulles. \emph{Renvoi : C4, definition.}
    \item \textbf{D} — La formule de l'union est P(A union B) = P(A) + P(B) - P(A inter B), pas un produit. \emph{Renvoi : C4, definition.}
  \end{itemize}
  \item \textbf{Q13}
  \begin{itemize}
    \item \textbf{A} — Incompatibles => P(A inter B) = 0, mais P(A)*P(B) > 0, donc pas independants. \emph{Renvoi : C4, contre-exemple.}
    \item \textbf{C} — On peut conclure : P(A inter B) = 0 != P(A)*P(B) > 0. \emph{Renvoi : C4, contre-exemple.}
    \item \textbf{D} — P\_A(B) = P(A inter B)/P(A) = 0/P(A) = 0, pas 1. \emph{Renvoi : C4, incompatibilite.}
  \end{itemize}
  \item \textbf{Q14}
  \begin{itemize}
    \item \textbf{A} — L'independance de A et B n'implique pas l'incompatibilite de Abar et B. \emph{Renvoi : C4, propriete.}
    \item \textbf{C} — Rien ne dit que P(Abar) = P(B). \emph{Renvoi : C4, propriete.}
    \item \textbf{D} — Il y a un lien precis : ils sont independants (propriete demontrable). \emph{Renvoi : C4, propriete.}
  \end{itemize}
  \item \textbf{Q15}
  \begin{itemize}
    \item \textbf{A} — P(M | T+) est la valeur predictive positive, pas la sensibilite. \emph{Renvoi : C5, definitions.}
    \item \textbf{C} — P(T- | Mbar) est la specificite, pas la sensibilite. \emph{Renvoi : C5, definitions.}
    \item \textbf{D} — P(M) est la prevalence, pas la sensibilite. \emph{Renvoi : C5, definitions.}
  \end{itemize}
  \item \textbf{Q16}
  \begin{itemize}
    \item \textbf{A} — P(T+ | M) est la sensibilite, pas la VPP. \emph{Renvoi : C5, definitions.}
    \item \textbf{C} — P(T+) est la probabilite globale d'un test positif. \emph{Renvoi : C5, definitions.}
    \item \textbf{D} — P(Mbar | T-) est la valeur predictive negative. \emph{Renvoi : C5, definitions.}
  \end{itemize}
  \item \textbf{Q17}
  \begin{itemize}
    \item \textbf{A} — L'eleve confond sensibilite et VPP. Ici $P(M\cap T^+)=0{,}00099$ et $P(T^+)=0{,}01098$, donc la VPP vaut environ $9\%$, pas $99\%$. \emph{Renvoi : C5, VPP et prevalence.}
    \item \textbf{B} — Avec une prevalence de $0{,}1\%$, les faux positifs representent $0{,}01\times0{,}999=0{,}00999$ de la population ; la VPP vaut environ $0{,}00099/0{,}01098=9\%$, pas $50\%$. \emph{Renvoi : C5, VPP et prevalence.}
    \item \textbf{D} — L'eleve recopie la prevalence. La VPP est la probabilite conditionnelle $P(M\mid T^+)$ et vaut ici environ $9\%$, pas $0{,}1\%$. \emph{Renvoi : C5, formule de Bayes.}
  \end{itemize}
  \item \textbf{Q18}
  \begin{itemize}
    \item \textbf{A} — Les taux de defaut dependent de la machine : l'origine et le defaut ne sont donc pas supposes independants. Il faut calculer la probabilite totale du defaut puis appliquer Bayes. \emph{Renvoi : C5, methode.}
    \item \textbf{C} — Les probabilites simples $P(M_i)$ et conditionnelles $P(D\mid M_i)$ doivent etre combinees pour calculer la probabilite inverse $P(M_i\mid D)$ ; elles ne suffisent pas prises separement. \emph{Renvoi : C5, methode.}
    \item \textbf{D} — Le TVI est un theoreme d'analyse, sans rapport avec les probabilites. \emph{Renvoi : C5, hors sujet.}
  \end{itemize}
\end{itemize}
\fi
% NEXUS-QCM-TEACHER-ONLY-END
